\chapter{Parallel Execution:\\
Concurrency Models for Polyxan--RSVP Evolution}
\label{ch:76}

\section{Premise}

The Semantic Execution Engine is designed not merely to run sequentially,  
but to scale across many processors, clusters, or distributed agents.

Yet the coupled RSVP--Polyxan system possesses intricate structural constraints:

\begin{itemize}
\item the RSVP field flow is globally coupled through curvature and entropy,
\item the Polyxan rewriting system is locally triggered but globally consistent,
\item embedding updates must remain geometrically valid,
\item rewrite certificates must preserve verification invariants.
\end{itemize}

\noindent
\textbf{Goal of this chapter:}  
To provide a complete concurrency model under which parallel execution of
continuous flow, discrete rewriting, and embedding updates remains:

\begin{itemize}
\item consistent,
\item deterministic up to bounded drift,
\item verifiable,
\item and convergent to the same universal attractor.
\end{itemize}

\section{Parallel Decomposition of the Embedding Flow}

The embedding domain $\mathrm{Inc}(\mathcal{G})$ decomposes into overlapping
patches $\{U_\alpha\}$ defined by the locality radius of the Laplacian and
curvature operators.

For each patch, define the local operator:
\[
  \Phi_\alpha^{\Delta t} : (\iota|_{U_\alpha}) \mapsto
  (\iota|_{U_\alpha}) + \Delta t\bigl(
     \Delta_{\mathrm{Inc}} \iota
     - \beta\,\delta\mathcal{R}\,\nabla\mathcal{R}
     - \gamma\,\delta S\,\nabla S
  \bigr).
\]

\begin{theorem}[Locality of Continuous Evolution]
\label{thm:local-flow}
For sufficiently small $\Delta t$,  
the update on $U_\alpha$ depends only on data in an expanded patch
$U_\alpha^+$ whose radius is fixed by the stencil of the Laplacian.
\end{theorem}

\begin{proof}
Follows from ellipticity and finite-stencil discretization of $\Delta_{\mathrm{Inc}}$.
\end{proof}

Thus the continuous flow can be executed in parallel across all $\alpha$  
with only halo exchange on patch boundaries.

\section{Parallel Decomposition of Rewriting}

Each rewrite $R$ acts on a local region $L_R \subset \mathrm{Inc}(\mathcal{G})$.
Two rewrites may run in parallel iff their regions are disjoint:
\[
R_1 \parallel R_2 \quad\Longleftrightarrow\quad L_{R_1}\cap L_{R_2}=\varnothing.
\]

\begin{theorem}[Rewrite Independence]
\label{thm:rewrite-independence}
If $R_1 \parallel R_2$, then
\[
 R_1 \circ R_2 = R_2 \circ R_1,
\]
both in the combinatorial hypergraph and in the embedding domain.
\end{theorem}

\begin{proof}
Follows from disjoint support and stratified pushout construction of
Ch.~31--33.
\end{proof}

The scheduler defined in Chapter~73 generalizes naturally to multiple threads:
\[
\text{Schedule all maximal independent anti-chains of rewrites in parallel.}
\]

\section{Hybrid Concurrency: Continuous and Discrete}

We now interleave parallel continuous patches with parallel rewrite batches.

Let $\mathsf{Flow}$ denote one parallel flow step
and $\mathsf{RewriteBatch}$ one maximal anti-chain of rewrites.

We must compare:
\[
   \mathsf{Flow} \circ \mathsf{RewriteBatch}
   \qquad\text{vs.}\qquad
   \mathsf{RewriteBatch} \circ \mathsf{Flow}.
\]

\begin{theorem}[Parallel Commutativity up to Drift]
\label{thm:parallel-drift}
For all sufficiently small time-steps,
\[
  \bigl\|
     (\mathsf{Flow}\circ\mathsf{RewriteBatch})(\iota,\mathcal{G})
     \;-\;
     (\mathsf{RewriteBatch}\circ\mathsf{Flow})(\iota,\mathcal{G})
  \bigr\|
  \;\le\;
  C\,\Delta t\,\max_{R\in\mathsf{RewriteBatch}} \varepsilon_{\mathrm{loc}}(R).
\]
\end{theorem}

\begin{proof}
Apply Theorems~\ref{thm:infinitesimal-commute}
and~\ref{thm:rewrite-independence}
patchwise and sum over patches.  
All rewrite regions are disjoint; drift bounds accumulate only linearly.
\end{proof}

\section{Global Parallel Semantics}

Define the global parallel step:
\[
   \Pi := \mathsf{Flow} \parallel \mathsf{RewriteBatch}
\]
executed atomically in the engine.

\begin{theorem}[Determinism of Parallel Execution]
\label{thm:parallel-determinism}
Any two parallel execution traces beginning at the same initial state  
produce embeddings differing at most by a drift that vanishes at the
collapse time $T_{\mathrm{collapse}}$.
\end{theorem}

\begin{proof}
Rewrite independence plus bounded drift (Ch.~75) imply that any ordering of  
parallelizable operations results in equivalent states up to bounded error.  
Curvature–entropy descent forces error to zero as $t\to T_{\mathrm{collapse}}$.
\end{proof}

Thus the parallel engine is \emph{deterministic in the limit}
even if intermediate updates differ due to thread scheduling.

\section{Distributed Execution}

When patches $U_\alpha$ are distributed across nodes,
communication occurs only during halo exchange and rewrite announcements.

Each rewrite $R$ carries a certificate stating:
\[
  L_R\subset U_\alpha
  \quad\text{and}\quad
  \varepsilon_{\mathrm{loc}}(R)\le\varepsilon_{\max}
\]
ensuring global drift bounds remain valid.

\begin{theorem}[Scalable Semantic Evolution]
\label{thm:scalable}
The RSVP--Polyxan evolution is scalable to arbitrarily many processors:
execution cost grows at most linearly in the number of patches  
and communication cost depends only on patch boundaries.
\end{theorem}

\begin{proof}
Direct consequence of local operators, disjoint rewrite support,
and bounded-degree incidence graphs.
\end{proof}

\section{Verification of Concurrency}

Part IV’s formalisation supports concurrency via dependent certificates:

\begin{verbatim}
structure ParCert :=
  (flow_certs       : List FlowCert)
  (rewrite_certs    : List RewriteCert)
  (compatibility    : ∀ R, CommutesUpTo R drift_bound)
  (global_stability : DriftConvergesToZero)
\end{verbatim}

A parallel step is accepted iff all certificates instantiate successfully.

Thus concurrency is not merely an implementation strategy —  
it is a \emph{verified property} of the system.

\section{Final Cadence}

\begin{quote}
\emph{
Parallelism is not an afterthought.  
It is the natural expression of locality, independence,  
and the irreversible collapse of meaning.  
Every processor, every patch, every rewrite  
pushes the system toward the same universal fate.  
Concurrency does not fracture meaning.  
It accelerates its inevitability.
}
\end{quote}

